\documentclass[11pt,letterpaper]{article}

\usepackage[T1]{fontenc}
\usepackage[margin=1in]{geometry}
\usepackage{booktabs}
\usepackage{listings}
\usepackage{xcolor}
\usepackage{hyperref}

\hypersetup{
  colorlinks=true,
  linkcolor=blue!70!black,
  urlcolor=blue!70!black,
  pdftitle={ebnf2tikz User Manual},
  pdfauthor={Larry D. Pyeatt}
}

\lstdefinelanguage{ebnf}{
  morekeywords={},
  morestring=[b]',
  morestring=[b]",
  morecomment=[l]{//},
  morecomment=[l]{---},
  morecomment=[s]{(*}{*)},
  literate={|}{{\textbar}}1
           {<<--}{{\textless\textless--}}4
           {-->>}{{--\textgreater\textgreater}}4,
  sensitive=true
}

\lstset{
  basicstyle=\ttfamily\small,
  keywordstyle=\bfseries,
  stringstyle=\color{red!70!black},
  commentstyle=\color{green!50!black}\itshape,
  showstringspaces=false,
  frame=single,
  framerule=0.4pt,
  xleftmargin=1em,
  xrightmargin=1em,
  breaklines=true,
  aboveskip=1em,
  belowskip=1em
}

\newcommand{\opt}[1]{\texttt{#1}}
\newcommand{\prog}{\texttt{ebnf2tikz}}

\title{\textbf{ebnf2tikz User Manual}}
\author{Larry D. Pyeatt}
\date{\today}

\begin{document}

\maketitle
\tableofcontents
\newpage

%% ----------------------------------------------------------------
\section{Introduction}

\prog{} is an optimizing compiler that converts Extended Backus-Naur
Form (EBNF) grammar specifications into railroad (syntax) diagrams
expressed as \LaTeX{} TikZ commands.  It supports annotations that
control inlining of productions, figure captions, and rotated figures.
An optimizer simplifies the resulting abstract syntax tree before
layout, producing compact, readable diagrams.

\prog{} is free software distributed under the GNU General Public
License, version~3 or later.

%% ----------------------------------------------------------------
\section{Installation}

\subsection{Prerequisites}

\begin{itemize}
  \item A C++ compiler with C++17 support (e.g., \texttt{g++})
  \item \href{https://www.gnu.org/software/bison/}{Bison} $\geq$ 3.7
  \item \href{https://github.com/westes/flex}{Flex}
  \item \href{https://cmake.org/}{CMake} $\geq$ 3.16
  \item \texttt{pdflatex} with TikZ (for rendering diagrams)
\end{itemize}

\subsection{Building}

\begin{lstlisting}[language=bash]
cmake .
make
\end{lstlisting}

This produces the \prog{} executable in the current directory.  To
verify the build:

\begin{lstlisting}[language=bash]
make check        # Run AST-dump regression tests
make check-tikz   # Build PDF with all railroad diagrams
\end{lstlisting}

%% ----------------------------------------------------------------
\section{Quick Start}

\begin{enumerate}
\item Write an EBNF grammar file, for example \texttt{mygrammar.ebnf}:
\begin{lstlisting}[language=ebnf]
greeting = 'hello' , [ name ] , '!' ;
name = identifier ;
\end{lstlisting}

\item Run \prog{} to produce initial TikZ output (with estimated node sizes):
\begin{lstlisting}[language=bash]
./ebnf2tikz mygrammar.ebnf mygrammar.tex
\end{lstlisting}

\item Include the output in a \LaTeX{} document with the required
  preamble (see Section~\ref{sec:latex}) and compile:
\begin{lstlisting}[language=bash]
pdflatex mydocument
\end{lstlisting}

\item Run \prog{} again to use the real node sizes written by
  \LaTeX{} to \texttt{bnfnodes.dat}:
\begin{lstlisting}[language=bash]
./ebnf2tikz mygrammar.ebnf mygrammar.tex
\end{lstlisting}

\item Compile \LaTeX{} again for the final output:
\begin{lstlisting}[language=bash]
pdflatex mydocument
\end{lstlisting}
\end{enumerate}

This multi-pass workflow is similar to the feedback loop required by
Bib\TeX{} or \texttt{makeindex}.  On the first pass, \prog{} uses
estimated node dimensions.  \LaTeX{} then measures the actual rendered
sizes of each node and writes them to \texttt{bnfnodes.dat} via the
\verb|\writenodesize| command.  On the second pass, \prog{} reads
these measurements and produces an accurate layout.

%% ----------------------------------------------------------------
\section{Command-Line Options}

\begin{center}
\begin{tabular}{@{}lll@{}}
\toprule
Short & Long & Description \\
\midrule
\opt{-O} & \opt{-{}-optimize}     & Enable optimization (default). \\
\opt{-n} & \opt{-{}-nooptimize}   & Disable optimization and subsumption. \\
\opt{-f} & \opt{-{}-makefigures}  & Wrap \texttt{tikzpicture}s in \LaTeX{} \texttt{figure} environments. \\
\opt{-d} & \opt{-{}-dumponly}     & Parse and dump the AST; do not generate TikZ output. \\
\opt{-B} & \opt{-{}-dump-before}  & Dump the AST before optimization (to stdout). \\
\opt{-A} & \opt{-{}-dump-after}   & Dump the AST after optimization (to stdout). \\
\opt{-h} & \opt{-{}-help}         & Print usage information and exit. \\
\bottomrule
\end{tabular}
\end{center}

\medskip\noindent
Usage:
\begin{lstlisting}[language=bash]
./ebnf2tikz [options] input_file output_file
\end{lstlisting}

The \opt{-O} flag is the default and is accepted for compatibility.
Use \opt{-n} to suppress both the optimizer and annotation-driven
subsumption.

The \opt{-B} and \opt{-A} flags print the AST to standard output
before and after optimization, respectively.  These are useful for
debugging grammars.  The \opt{-d} flag prints the AST and skips TikZ
generation entirely.

%% ----------------------------------------------------------------
\section{EBNF Syntax Reference}

\prog{} accepts standard EBNF as defined by ISO~14977, with a few
extensions.

\subsection{Productions}

A grammar consists of one or more productions.  Each production has a
name, an equals sign, a body expression, and a terminating semicolon:

\begin{lstlisting}[language=ebnf]
production_name = body_expression ;
\end{lstlisting}

\subsection{Constructs}

\begin{center}
\begin{tabular}{@{}llp{6cm}@{}}
\toprule
Syntax & Meaning & Example \\
\midrule
\texttt{A , B}     & Sequence (concatenation) & \texttt{'if' , expr , 'then'} \\
\texttt{A | B}     & Choice (alternation)     & \texttt{'true' | 'false'} \\
\texttt{[ A ]}     & Optional                 & \texttt{[ 'else' , statement ]} \\
\texttt{\{ A \}}   & Repetition (zero or more) & \texttt{\{ ',' , param \}} \\
\texttt{( A )}     & Grouping                 & \texttt{( 'a' | 'b' ) , 'c'} \\
\texttt{'text'}    & Terminal (single quotes)  & \texttt{'while'} \\
\texttt{"text"}    & Terminal (double quotes)  & \texttt{"begin"} \\
\texttt{? text ?}  & Special sequence          & \texttt{? any character ?} \\
\texttt{name}      & Nonterminal reference     & \texttt{expression} \\
\texttt{\textbackslash\textbackslash} & Explicit line break & Forces a new row in the diagram \\
\bottomrule
\end{tabular}
\end{center}

\subsection{Comments}

Three comment styles are supported:

\begin{lstlisting}[language=ebnf]
// C++-style line comment
--- Ada-style line comment
(* Block comment, can span
   multiple lines *)
\end{lstlisting}

%% ----------------------------------------------------------------
\section{Annotations}
\label{sec:annotations}

Annotations are placed immediately before a production definition
using the \verb|<<-- ... -->>| syntax.  Multiple annotations may be
placed before a single production.

\subsection{\texttt{subsume}}

\begin{lstlisting}[language=ebnf]
<<-- subsume -->>
production_name = body ;
\end{lstlisting}

The \texttt{subsume} annotation marks a production for inlining:
wherever its name appears as a nonterminal reference in other
productions, the reference is replaced by the production's body.  The
subsumed production itself is not emitted as a separate diagram.

\subsection{\texttt{subsume as}}

\begin{lstlisting}[language=ebnf]
<<-- subsume as regex_pattern -->>
production_name = body ;
\end{lstlisting}

Like \texttt{subsume}, but the production is only inlined into
nonterminals whose names match the given POSIX extended regular
expression.

\subsection{\texttt{caption}}

\begin{lstlisting}[language=ebnf]
<<-- caption A descriptive caption -->>
production_name = body ;
\end{lstlisting}

Sets the \LaTeX{} figure caption for this production's diagram.
Requires the \opt{-f} flag to have any effect, since captions only
apply to \texttt{figure} environments.

\subsection{\texttt{sideways}}

\begin{lstlisting}[language=ebnf]
<<-- sideways -->>
production_name = body ;
\end{lstlisting}

Places this production's diagram in a \texttt{sidewaysfigure}
environment (from the \texttt{rotating} package) instead of a regular
\texttt{figure}.  Requires the \opt{-f} flag.

%% ----------------------------------------------------------------
\section{\LaTeX{} Integration}
\label{sec:latex}

There is currently no \texttt{ebnf2tikz.sty} package file.  To use
\prog{} output in your \LaTeX{} document, you need to define the
required packages, lengths, TikZ styles, and the
\verb|\writenodesize| command in your preamble.

\subsection{Required Packages}

\begin{lstlisting}[language={[LaTeX]TeX}]
\usepackage{tikz,pgf}
\usetikzlibrary{positioning,shapes,matrix,chains,
  scopes,fit,calc,shadings,intersections}
\usepackage[nomessages]{fp}
\usepackage{rotating}     % only if using sideways annotations
\end{lstlisting}

\subsection{Font Commands}

\begin{lstlisting}[language={[LaTeX]TeX}]
\newcommand{\railname}[1]{\it #1}
\newcommand{\railtermname}[1]{\bf #1}
\end{lstlisting}

The \verb|\railname| command formats nonterminal labels and
\verb|\railtermname| formats terminal labels in the production title
area.

\subsection{Lengths}

\begin{lstlisting}[language={[LaTeX]TeX}]
\newlength\railcolsep
\setlength\railcolsep{5pt}
\newlength\railrowsep
\setlength\railrowsep{\baselineskip}
\newlength\railnodeheight
\setlength\railnodeheight{14pt}
\newlength\railcorners
\setlength\railcorners{\railcolsep}
\addtolength\railcorners{-1pt}
\end{lstlisting}

\begin{center}
\begin{tabular}{@{}lp{8cm}@{}}
\toprule
Length & Purpose \\
\midrule
\verb|\railcolsep|     & Horizontal gap between nodes. \\
\verb|\railrowsep|     & Vertical gap between rows in multi-row diagrams. \\
\verb|\railnodeheight| & Minimum height of terminal and nonterminal nodes. \\
\verb|\railcorners|    & Rounding radius for rail corners. \\
\bottomrule
\end{tabular}
\end{center}

\subsection{TikZ Styles}

\begin{lstlisting}[language={[LaTeX]TeX}]
\tikzset{
  nonterminal/.style={
    draw, inner sep=2pt, rectangle,
    minimum size=\railnodeheight, thick,
    top color=white,
    bottom color=red!50!black!20,
    font=\itshape
  },
  terminal/.style={
    draw, inner sep=2pt, rounded rectangle,
    minimum size=\railnodeheight, thick,
    top color=white, bottom color=black!20,
    font=\ttfamily
  },
  every picture/.style={thick},
  every node/.style={inner sep=0, outer sep=0, draw=none}
}
\end{lstlisting}

The \texttt{nonterminal} style renders rectangular nodes with a
gradient fill.  The \texttt{terminal} style renders rounded rectangles.
You may customize these styles to suit your document's appearance.

\subsection{The \texttt{\textbackslash writenodesize} Command}

This command is essential for the multi-pass workflow.  It measures a
TikZ node and writes its dimensions to \texttt{bnfnodes.dat}:

\begin{lstlisting}[language={[LaTeX]TeX}]
\newwrite\tempfile
\immediate\openout\tempfile=bnfnodes.dat

\makeatletter
\newcommand{\writenodesize}[1]{%
  \pgfpointdiff{\pgfpointanchor{#1}{west}}%
               {\pgfpointanchor{#1}{east}}%
  \FP@eval\@temp@x{\pgfmath@tonumber{\pgf@x}}%
  \pgfpointdiff{\pgfpointanchor{#1}{south}}%
               {\pgfpointanchor{#1}{north}}%
  \FP@eval\@temp@y{\pgfmath@tonumber{\pgf@y}}%
  \immediate\write\tempfile{#1 \@temp@x,  \@temp@y}}
\makeatother
\end{lstlisting}

You must also write the \verb|\textwidth| value to \texttt{bnfnodes.dat}
at the beginning of the document body so that \prog{} knows the
available width for auto-wrapping:

\begin{lstlisting}[language={[LaTeX]TeX}]
\begin{document}
\makeatletter
\pgf@x=\textwidth%
\FP@eval\@temp@tw{\pgfmath@tonumber{\pgf@x}}%
\immediate\write\tempfile{textwidth \@temp@tw}
\makeatother
\end{lstlisting}

At the end of the document, close the file:
\begin{lstlisting}[language={[LaTeX]TeX}]
\immediate\closeout\tempfile
\end{document}
\end{lstlisting}

\subsection{Scratch Lengths}

The generated TikZ code uses three scratch lengths.  Declare them in
your preamble or document body before including \prog{} output:

\begin{lstlisting}[language={[LaTeX]TeX}]
\newlength{\tmplength}
\newlength{\tmplengthb}
\newlength{\tmplengthc}
\setlength{\tmplength}{12pt}
\setlength{\tmplengthb}{12pt}
\setlength{\tmplengthc}{12pt}
\end{lstlisting}

%% ----------------------------------------------------------------
\section{Auto-Wrap}
\label{sec:autowrap}

Productions whose body exceeds the available \verb|\textwidth| are
automatically broken across multiple rows in the railroad diagram.
This happens without any user intervention.

Additionally, wide nonterminal labels that contain spaces are
automatically wrapped using \verb|\shortstack| to reduce their
rendered width.

If you need to force a line break at a specific point, use the
explicit line break operator \verb|\\| in your EBNF input:

\begin{lstlisting}[language=ebnf]
long_production =
    first_part , second_part , \\ third_part , fourth_part ;
\end{lstlisting}

%% ----------------------------------------------------------------
\section{Examples}
\label{sec:examples}

\subsection{Case Statement Alternative}

This example from the VHDL grammar demonstrates subsumption.  Three
productions are defined, but \texttt{choices} and \texttt{choice} are
annotated with \texttt{subsume}, so they are inlined into
\texttt{case\_statement\_alternative}:

\begin{lstlisting}[language=ebnf]
case_statement_alternative =
    'when' , choices , '=>', sequence_of_statements;

<<-- subsume -->>
choices =
    choice, { '|', choice } ;

<<-- subsume -->>
choice =
    simple_expression |
    discrete_range |
    simple_name |
    'others' ;
\end{lstlisting}

With optimization enabled (the default), the optimizer inlines the
subsumed productions and simplifies the resulting AST.  The final
diagram shows a single, compact railroad for
\texttt{case\_statement\_alternative} with all the choices expanded
inline.

\subsection{Tool List}

This example shows a more complex grammar with subsumption:

\begin{lstlisting}[language=ebnf]
<<-- subsume -->>
toollist =
  '(' , tool , { ',' , tool } , ')' ;

<<-- subsume -->>
tool =
  'hammer' | 'screwdriver' | saw;

<<-- subsume -->>
saw =
  ('hand' | ( [ 'cordless'] ,
  ('reciprocating' | 'circular') ,
  'power' ) ),'saw';
\end{lstlisting}

When all three productions are marked with \texttt{subsume}, each
nonterminal reference is replaced by the referenced production's body.
The result is a single, fully-expanded railroad diagram.  In practice,
you would mark only the helper productions for subsumption, leaving
the top-level production (\texttt{toollist}) as a standalone diagram
that incorporates the others.

%% ----------------------------------------------------------------
\section{Tips}
\label{sec:tips}

\begin{itemize}
\item Use the \opt{-B} and \opt{-A} flags to inspect the AST before
  and after optimization.  This is helpful for understanding how the
  optimizer transforms your grammar.

\item When adding \prog{} diagrams to an existing document, it is
  convenient to \verb|\input| the generated TikZ file rather than
  pasting it inline.

\item The \opt{-f} flag wraps each diagram in a \texttt{figure}
  environment, which enables captions, labels, and
  \verb|\listoffigures|.  Without \opt{-f}, diagrams are placed
  inline.

\item Production names must be valid identifiers (letters, digits,
  underscores).  Underscores in production names are rendered in the
  diagram using \LaTeX{} formatting.

\item The feedback loop (\prog{} $\to$ \texttt{pdflatex} $\to$
  \prog{} $\to$ \texttt{pdflatex}) need only be repeated when
  node content changes.  If you are only editing prose around
  existing diagrams, a single \texttt{pdflatex} run suffices.
\end{itemize}

\end{document}
